\documentclass[a4paper]{article}
% Español (México) con XeLaTeX: fontspec para fuentes Unicode, polyglossia
% para la división silábica y los nombres del documento en la variante mexicana.
\usepackage{fontspec}
\usepackage{polyglossia}
\setdefaultlanguage[variant=mexican]{spanish}
% Los nombres chinos van como «pinyin (original)»: xeCJK da a los caracteres
% Han su propia fuente; sin ella XeLaTeX los omite con un aviso.
% FandolSong no cubre algunos caracteres de nombres (U+73FA); WenQuanYi Zen Hei sí.
\usepackage[AutoFallBack=true]{xeCJK}
\setCJKmainfont{FandolSong-Regular.otf}
\setCJKfallbackfamilyfont{\CJKrmdefault}{WenQuanYi Zen Hei}
% polyglossia no traduce los nombres de listings.
\AtBeginDocument{\providecommand{\lstlistingname}{}\renewcommand{\lstlistingname}{Listado}%
  \providecommand{\lstlistlistingname}{}\renewcommand{\lstlistlistingname}{Índice de listados}}
\usepackage{amsmath, amssymb}
\usepackage{graphicx}
\usepackage[margin=2.35cm]{geometry}
\usepackage[most]{tcolorbox}
\usepackage{booktabs}
\usepackage{tabularx}
\usepackage{array}
\usepackage{float}
\usepackage{xcolor}
\usepackage{hyperref}
\usepackage{fancyhdr}
\setCJKmainfont[AutoFakeBold=2.4]{AR PL UMing CN}
\setCJKsansfont[AutoFakeBold=2.4]{Droid Sans Fallback}
\setCJKmonofont[AutoFakeBold=2.4]{Droid Sans Fallback}
\hypersetup{colorlinks=true, linkcolor=blue!55!black, urlcolor=blue!60!black}
\graphicspath{{./}{figures/}}
\setlength{\parskip}{0.55em}
\setlength{\parindent}{2em}
\setlength{\emergencystretch}{3em}
\renewcommand{\arraystretch}{1.18}
\newtcolorbox{knowledgebox}[1]{enhanced, colback=blue!5!white, colframe=blue!70!black, colbacktitle=blue!70!black, coltitle=white, fonttitle=\bfseries, title={#1}, sharp corners, breakable}
\newtcolorbox{importantbox}[1]{enhanced, colback=yellow!10!white, colframe=yellow!75!black, colbacktitle=yellow!75!black, coltitle=black, fonttitle=\bfseries, title={#1}, sharp corners, breakable}
\newtcolorbox{warningbox}[1]{enhanced, colback=red!5!white, colframe=red!70!black, colbacktitle=red!70!black, coltitle=white, fonttitle=\bfseries, title={#1}, sharp corners, breakable}
\pagestyle{fancy}
\fancyhf{}
\fancyfoot[C]{\thepage}
\fancyfoot[R]{\footnotesize\color{gray}\href{https://github.com/hqhq1025/ai-course-notes}{AI Course Notes}}
\renewcommand{\headrulewidth}{0pt}
\renewcommand{\footrulewidth}{0pt}
\newcommand{\videourl}{https://www.youtube.com/watch?v=a04POJEknCY}
\newcommand{\repourl}{https://github.com/hqhq1025/ai-course-notes}

\begin{document}
\begin{titlepage}
\centering
{\Large Notas didácticas de una entrevista a fondo\par}
\vspace{1.0cm}
{\huge\bfseries YouWare, OS Agent y el nuevo paradigma de producto en el emprendimiento de aplicaciones de IA}\par
\vspace{0.5cm}
{\Large Zhang Xiaojun conversa con Ming Chaoping (明超平), fundador de YouWare\par}
\vspace{0.35cm}
{\large Elaborado a partir del video público, la transcripción, la estructura de capítulos, la descripción del video y diagramas conceptuales propios\par}
\vspace{0.3cm}
{\large 2025-05-29\par}
\vspace{0.85cm}
\includegraphics[width=0.88\textwidth,height=0.42\textheight,keepaspectratio]{cover.jpg}\par
\vfill
\begin{tcolorbox}[width=0.92\textwidth, colback=black!2!white, colframe=black!55, sharp corners]
\textbf{Título del video}: 101. Entrevista de 3 horas con Ming Chaoping, fundador de YouWare: hoy el Agent es como un gorila que acaba de tomar un atizador\par
\textbf{Canal del video}: Zhang Xiaojun Podcast\par
\textbf{Duración del video}: 02:42:45\par
\textbf{Enlace del video}: \href{\videourl}{\nolinkurl{\videourl}}\par
\textbf{Nota sobre los materiales}: YouTube no ofrece subtítulos descargables; el texto se elaboró a partir de una transcripción por bloques con faster-whisper large-v3 (CPU, int8), la descripción del video, la estructura de capítulos y figuras didácticas propias. Las tomas fijas de la entrevista no se reutilizan en el cuerpo del texto.\par
\textbf{Página del proyecto}: \href{\repourl}{\nolinkurl{\repourl}}\par
\end{tcolorbox}
\end{titlepage}

\begingroup
\small
\setlength{\parskip}{0pt}
\renewcommand{\baselinestretch}{0.92}\selectfont
\tableofcontents
\endgroup
\newpage

\section{Introducción: en 2025 el protagonismo pasa de las empresas de modelos a las aplicaciones y los Agent}

Esta sección sitúa el episodio. Zhang Xiaojun (张小珺) explica al inicio que en 2024 el relato de la IA todavía giraba en torno a las empresas de modelos grandes, y que en 2025 los nuevos protagonistas son las empresas de aplicaciones y los Agent. Ming Chaoping (明超平) es un emprendedor de aplicaciones de IA nacido en 1995; trabajó en producto sucesivamente en OnePlus, ByteDance y Moonshot, es CEO por primera vez y el primer producto que lanzó se llama YouWare. El valor de este episodio no está solo en escuchar la historia personal de un founder joven, sino en ver cómo un fundador con perfil de producto enlaza teléfonos móviles, video corto, Kimi, coding, Agent y la oportunidad del OS en un mismo conjunto de juicios.

\begin{importantbox}{Tesis central del episodio}
La oportunidad de YouWare no se reduce a «hacer un producto de Agent»: consiste en redefinir el contenedor de creación y el entorno de ejecución de la era de la IA. El juicio de Ming Chaoping es que los modelos base seguirán mejorando, pero las empresas de aplicaciones deben adaptar modelos más inteligentes a necesidades de producción concretas mediante el entorno, la experiencia, la interacción y los efectos de red.
\end{importantbox}

\begin{knowledgebox}{Nota sobre la estrategia visual}
Este video es una entrevista con encuadre fijo, sin slides, pizarrón ni demostración de producto. El texto usa la portada solo para identificar la fuente; todas las imágenes del cuerpo son diagramas conceptuales de elaboración propia, que explican las tres etapas como gerente de producto, experiencia y datos, la bitter lesson, el consumo de token, el contenedor/entorno, el Vibe Coder, AgentRank, el OS Agent y la confrontación con el Ego.
\end{knowledgebox}

\begin{figure}[H]
\centering
\includegraphics[width=0.96\textwidth]{product-manager-three-stations.png}
\caption{Tres etapas como gerente de producto: OnePlus, ByteDance y Moonshot forman la base de su juicio de producto. Diagrama conceptual de elaboración propia, organizado a partir de la conversación de 00:03:16--01:05:05.}
\end{figure}

\begin{knowledgebox}{Lectura de la figura: cada una de las tres experiencias entrena una capacidad de producto distinta}
OnePlus entrena la experiencia y la sensibilidad estética; ByteDance, los datos, el crecimiento y la iteración de alta frecuencia; Moonshot, los modelos, los productos de IA y la incertidumbre de la frontera. YouWare no surge de la nada: es la recombinación de estas tres capacidades en la ventana de oportunidad de las aplicaciones de Agent.
\end{knowledgebox}

\subsection{Por qué esta no es una entrevista de emprendimiento común}

Esta sección desplaza primero la atención de la trayectoria personal al método de producto. La historia de Ming Chaoping tiene dos líneas: una es cómo una persona forma su juicio de producto; la otra, cómo una empresa de aplicaciones de IA encuentra su lugar junto a los gigantes de los modelos. La primera incluye el entrenamiento en debate, la competencia de vehículos inteligentes, el entrenamiento en experiencia de OnePlus y el entrenamiento en datos de ByteDance; la segunda incluye la bitter lesson, el consumo de token, el entorno como contenedor, AgentRank y el OS Agent.

\begin{knowledgebox}{Nota de clase: leer la historia personal como metodología}
La lectura útil de este episodio no es «el founder joven tiene mucha personalidad», sino «qué experiencias lo llevaron a qué juicios de producto». El debate corresponde a la perspectiva de un tercero; la competencia de vehículos inteligentes, al reward y a la retroalimentación práctica; OnePlus, a la experiencia; ByteDance, a los datos; Moonshot, al sentido de los límites en la era de los modelos.
\end{knowledgebox}

\subsection{Ruta de lectura}

Esta sección propone cómo leer el episodio. La primera mitad de la entrevista trata del debate, OnePlus, ByteDance y Kimi, y es fácil malinterpretarla como un currículum personal; la segunda mitad trata de los token, el contenedor, el Vibe Coder, AgentRank y el OS Agent, que son el núcleo del juicio de producto. La lectura correcta es: la primera mitad explica por qué hace producto de esta manera; la segunda, por qué las empresas de aplicaciones de IA todavía tienen oportunidad.

\noindent\begin{tabularx}{\textwidth}{p{0.18\textwidth}p{0.34\textwidth}X}
\toprule
Pregunta de lectura & Material de la entrevista & Juicio que se debe formar \\
\midrule
De dónde viene la intuición de producto & Debate, OnePlus, ByteDance, Kimi & La moldean en conjunto la perspectiva de un tercero, la sensibilidad a la experiencia, la infraestructura de datos y la ventana de los modelos. \\
Cómo evaluar las aplicaciones de IA & bitter lesson, consumo de token, per token valuation & No pulir adornos; reconstruir el producto en torno a la capacidad del modelo y a la eficiencia con que se consume la inteligencia. \\
De dónde surge el ecosistema de Agent & AgentRank, efectos de red, documento de identidad, confianza & Los Agent no son herramientas aisladas: forman redes de invocación, clasificación y confianza. \\
Cómo construir los mecanismos del CEO & Valor emocional, andar como sobre hielo delgado, confrontar el Ego, mentalidad de jugador de ajedrez & Emprender en aplicaciones exige que el fundador se recalibre continuamente a sí mismo y a su organización. \\
\bottomrule
\end{tabularx}

\subsection{Resumen de la sección}

La línea principal de EP101 es que el emprendimiento en aplicaciones de IA pasa del culto a los modelos a la construcción de entornos. El caso de YouWare muestra que las empresas de aplicaciones no pueden limitarse a esperar a que los modelos se vuelvan más potentes: deben decidir en qué contenedor se coloca la inteligencia, cómo se consume y cómo genera efectos de red.
\section{La base del gerente de producto: perspectiva de tercero, experiencia y datos}

La sección anterior explicó que este episodio debe leerse en clave de juicio de producto; esta sección examina primero de dónde proviene esa capacidad de juicio. En la universidad, Ming Chaoping (明超平) participaba en debates y destacaba sobre todo como cuarto orador; concluyó que lo más importante en un debate no es convencer al adversario, sino conservar siempre la perspectiva de un tercero. Ese entrenamiento se convirtió después en un método de producto: el usuario no es el adversario de las discusiones internas del equipo, sino el jurado que observa desde el público. En el instante en que abre el producto, el usuario no conoce tu historia, tu estrategia ni tu esfuerzo; solo sabe si puede entenderlo y usarlo de inmediato.

\begin{importantbox}{Nota de clase: volverse ingenuo en un segundo}
Ming Chaoping cita la intuición de producto al estilo de Zhang Xiaolong (张小龙): hay que convertirse en un usuario común que entra en el escenario por primera vez. Cualquier software que requiera explicaciones largas, ventanas emergentes y una instrucción compleja para poder usarse no ha completado de verdad su diseño de experiencia.
\end{importantbox}

\subsection{La experiencia no son los datos}

Después de la perspectiva de tercero, esta subsección aborda la primera lección que le dio OnePlus: un producto no es un conjunto de indicadores, sino la sensación continua del usuario en una situación concreta. La pregunta que hay que responder aquí es por qué una función que los datos parecen respaldar puede, aun así, generar ansiedad en el usuario. La autonomía de un teléfono tiene datos objetivos, pero la experiencia de autonomía no equivale a la curva de carga de la batería. Lo que más angustia al usuario puede ser la espera de carga entre 95--100\%, o el estado de batería casi agotada entre 0--5\%; una autonomía brillante en el papel no significa una buena experiencia. Este caso muestra que los datos pueden decirte qué ocurrió, pero no te dicen automáticamente cómo lo percibe el usuario.

\begin{figure}[H]
\centering
\includegraphics[width=0.94\textwidth]{experience-not-data.png}
\caption{La experiencia no son los datos: los datos de autonomía y la experiencia de autonomía no son equivalentes. Diagrama conceptual propio, elaborado a partir de la conversación de 00:26:37--00:33:58.}
\end{figure}

\begin{knowledgebox}{Lectura de la figura: la experiencia combina datos y umbrales psicológicos}
A la izquierda, Data son indicadores como autonomía promedio, retención y clics; a la derecha, Experience son la ansiedad, los puntos sensibles y la percepción subjetiva. Al leer esta figura conviene notar que la experiencia no se opone a los datos, sino que exige saber qué intervalos de los datos son psicológicamente más sensibles para el usuario.
\end{knowledgebox}

\begin{knowledgebox}{Términos clave: datos, experiencia, retrovisor}
\noindent\begin{tabularx}{\textwidth}{p{0.18\textwidth}p{0.36\textwidth}X}
\toprule
Concepto & Significado & Juicio en este episodio \\
\midrule
Datos & Registro cuantificado del comportamiento pasado & Permiten verificar resultados, pero no generan una dirección por sí mismos. \\
Experiencia & Percepción subjetiva del usuario en umbrales concretos & Determina si el producto se entiende, genera confianza y se sigue usando. \\
Retrovisor & Los datos son como el espejo retrovisor al manejar & Ayudan a evaluar el trayecto recorrido, pero no pueden ver por ti el camino de enfrente. \\
Juicio de convicción & Los datos no lo respaldan a corto plazo, pero la dirección sigue mereciendo la apuesta & Es especialmente importante en la ventana de las aplicaciones de IA. \\
\bottomrule
\end{tabularx}
\end{knowledgebox}

\subsection{Las fortalezas de ByteDance y sus efectos secundarios}

Al pasar del entrenamiento en experiencia de OnePlus a ByteDance (字节), el problema se convierte en otro régimen de velocidad y escala. Esta subsección examina la segunda lección que le dio ByteDance: la toma de decisiones guiada por datos y la iteración de alta frecuencia. OnePlus lanza un teléfono al año; ByteDance puede publicar una versión cada dos semanas, y el gerente de producto debe moverse con rapidez entre datos, experimentos, crecimiento y coordinación organizacional. Ming Chaoping reconoce que ByteDance le permitió dominar métodos de producto a mayor escala, pero también señala los efectos secundarios: un sistema de datos muy fuerte puede sofocar las ideas creativas que surgen de un destello de inspiración, porque algo verdaderamente nuevo puede mostrar datos negativos a corto plazo.

\begin{warningbox}{La infraestructura de datos no sustituye el juicio creativo}
Si solo se persigue un resultado positivo en pruebas A/B de corto plazo, el equipo será cada vez más hábil para optimizar partes existentes, pero no necesariamente creará formas nuevas. Esto vale en especial para el emprendimiento de aplicaciones de IA: los usuarios nuevos, los grupos de personas nuevos y los flujos de trabajo nuevos rara vez tienen datos estables al principio.
\end{warningbox}

\subsection{La competencia de autos inteligentes: un entrenamiento temprano de reward model}

Antes de su carrera como gerente de producto, Ming Chaoping menciona que en sus dos últimos años de universidad prácticamente dormía en el laboratorio mientras preparaba la competencia de autos inteligentes. Esa experiencia se entiende bien en el contexto del emprendimiento de aplicaciones de IA: que el auto vaya más rápido, que esquive obstáculos con más estabilidad o que la comunicación sea más fluida se refleja de inmediato en el resultado de la competencia. Él lo llama un muy buen reward model: los cambios en el código, los sensores, el hardware y la colaboración del equipo regresan al resultado de una forma muy directa.

Esta experiencia explica por qué más adelante fue sensible al «entorno» y a la «retroalimentación». Muchos gerentes de producto solo toman decisiones frente a pantallas y tableros de datos, mientras que la competencia de autos inteligentes lo puso desde temprano en contacto con ciclos cerrados del mundo real: entrada de sensores, salida de control, bugs, velocidad y resultados de la competencia. Aunque YouWare es un producto de software, también busca un ciclo cerrado similar: el usuario expresa una intención, el sistema genera una obra interactiva, la obra se comparte y recibe retroalimentación, y en la siguiente ronda el producto se ajusta.

\begin{importantbox}{Experiencia práctica: una buena retroalimentación de producto debe parecerse al cronómetro de una competencia}
Si la retroalimentación de un sistema es demasiado lenta o difusa, al equipo le cuesta formar intuición. La retroalimentación intensa de la competencia de autos inteligentes genera adicción, y las aplicaciones de IA necesitan un mecanismo similar: si el usuario publica, si comparte, si otros interactúan con lo que hizo, si sigue creando; todo eso se acerca más a un reward real que un clic aislado.
\end{importantbox}

\subsection{Resumen de la sección}

La base de producto de Ming Chaoping proviene de cuatro entrenamientos: el debate entrenó la perspectiva de tercero, la competencia de autos inteligentes entrenó la retroalimentación en ciclo cerrado, OnePlus entrenó la sensibilidad a la experiencia y ByteDance entrenó los datos y la iteración. En conjunto, estas capacidades explican por qué más adelante insistió en el contenedor, el entorno y la per token valuation.
\section{Bitter Lesson: no ir en contra de la mayor variable de la era de los modelos}

La sección anterior trató la base del producto; esta sección aborda el primer gran giro en el emprendimiento de aplicaciones de IA. Poco después de dejar Moonshot para emprender, Ming Chaoping (明超平) se topó con la bitter lesson: cuando la capacidad de los modelos es la mayor variable de la época, si una empresa de aplicaciones añade por instinto demasiado andamiaje, plantillas y ornamentos al producto, puede terminar alejándose de la dirección en la que avanzan los modelos. Detener el producto antes de lanzarlo no se debió a falta de trabajo, sino a que se dieron cuenta de que habían apostado por la variable equivocada.

\begin{figure}[H]
\centering
\includegraphics[width=0.96\textwidth]{bitter-lesson-pivot.png}
\caption{Giro de la Bitter Lesson: querer ornamentar y añadir andamiaje puede alejar el producto de la mayor variable de la época. Diagrama conceptual propio, elaborado a partir de la conversación de 01:05:05--01:13:18.}
\end{figure}

\begin{knowledgebox}{Lectura de la figura: saber detenerse también es una capacidad de producto}
La figura va de Scaffold a Launch?, Wrong variable, Stop y Rethink. Lo esencial aquí no es «no hacer nada», sino que la empresa de aplicaciones debe juzgar qué experiencias debe asumir el producto y qué capacidades conviene dejar que surjan de forma natural en el modelo. Confundir el andamiaje con una ventaja defensiva hace que el producto se vuelva cada vez más pesado.
\end{knowledgebox}

\subsection{El auge y la caída repentina de Noisee}

Esta subsección completa un giro importante de la etapa en Moonshot. En la entrevista se menciona el auge y la caída repentina de Noisee, un producto del extranjero, que le mostró a Ming Chaoping dos rasgos de la ventana de oportunidad de los productos de IA: el despegue puede llegar muy rápido, pero saber si el producto realmente sigue las capacidades del modelo y los escenarios de uso exige un juicio más sereno. Un producto puede recibir retroalimentación positiva de corto plazo en difusión, pero si la variable de fondo es incorrecta, seguir ornamentándolo solo aleja al equipo de la línea principal de la época.

Esta experiencia y la bitter lesson se explican mutuamente. Emprender en aplicaciones de IA no consiste en «hacer un demo vistoso y luego pagar por tráfico», sino en juzgar qué tipo de capacidad están liberando los modelos y en qué entorno los usuarios pueden convertir esa capacidad más rápido en obras o en productividad. Noisee le dio una percepción directa de la frontera; YouWare, en cambio, fue la elección de un nuevo punto de entrada, más cercano a la creación con coding y al entorno de publicación.

\begin{warningbox}{La retroalimentación positiva en difusión no significa que el rumbo del producto sea correcto}
Los productos de IA logran difusión con facilidad gracias a la novedad, pero la difusión no sustituye la retención, la creación, la reutilización ni el ecosistema. Los giros de Ming Chaoping en Moonshot y en las primeras etapas de su emprendimiento consistieron precisamente en distinguir entre «volverse popular» y «tener el rumbo correcto».
\end{warningbox}

\subsection{Velocidad de consumo de tokens y per token valuation}

La revelación que tuvo tras una noche de insomnio fue que uno de los indicadores clave de la era de la IA podría ser la velocidad de consumo de tokens. Los usuarios no gastan tokens por «gastar tokens», sino para que la inteligencia se convierta más rápido en resultados. Un buen escenario de Agent debería completar tareas más grandes con menos tokens, o hacer que los tokens se conviertan con mayor eficiencia en obras, sitios web, juegos, procesos y productividad.

\begin{figure}[H]
\centering
\includegraphics[width=0.96\textwidth]{token-consumption-speed.png}
\caption{Velocidad de consumo de tokens: en la era de la IA, los indicadores de producto pasan de clics/tiempo de uso a un consumo más eficiente de inteligencia. Diagrama conceptual propio, elaborado a partir de la conversación de 01:13:18--01:16:33 y 02:03:33--02:14:50.}
\end{figure}

\begin{knowledgebox}{Lectura de la figura: el token es consumo de inteligencia, no solo un costo}
En la figura, User intent entra al Container, luego Agent action consume Tokens y finalmente estos se convierten en Value. Al leer la figura, no hay que ver el token solo como costo de API; para una empresa de aplicaciones, el token es la unidad con la que se moviliza la inteligencia, y lo decisivo es cuánto valor para el usuario crea cada token.
\end{knowledgebox}

\[
\mathrm{PTV}=\frac{\mathrm{User\ Value}}{\mathrm{Tokens\ Consumed}}
\]

Donde $\mathrm{PTV}$ denota la per token valuation; $\mathrm{User\ Value}$ es la obra, la eficiencia, el ingreso o la experiencia que el usuario obtiene realmente; y $\mathrm{Tokens\ Consumed}$ son los tokens del modelo consumidos para completar la tarea. Esta expresión no es una fórmula financiera, sino una herramienta para pensar el producto: si se consumen muchos tokens y solo se produce una conversación casual, la densidad de valor es baja; si con pocos tokens se genera un sitio web o un juego que se puede compartir, con el que se puede interactuar y que se puede seguir editando, la densidad de valor es alta.

\begin{warningbox}{No interpretar el consumo de tokens como «cuanto más, mejor»}
La velocidad de consumo importa, pero no debe entenderse burdamente como quemar tokens. El objetivo real es convertir la inteligencia, de forma más rápida y eficiente, en valor perceptible para el usuario. Las conversaciones largas sin propósito, el ensayo y error repetido y el razonamiento en vacío solo elevan el costo y no forman una ventaja defensiva para el producto.
\end{warningbox}

\subsection{Resumen de la sección}

La primera lección del emprendimiento en aplicaciones de IA es seguir la mayor variable de la era de los modelos, en lugar de envolver al modelo en una capa compleja por hábitos de productos anteriores. Los indicadores clave de YouWare pasan de los clics/tiempo de uso tradicionales a cómo la inteligencia se encapsula en un contenedor, se consume con eficiencia y se convierte en obras.
\section{Contenedor y entorno: ofrecer solo el cuadro de entrada de un Chatbot es una irresponsabilidad}

La sección anterior trató sobre los tokens; esta trata sobre la forma de producto que los aloja. Ming Chaoping (明超平) considera que la palabra «cascarón» subestima el valor de las empresas de aplicaciones; lo más preciso es hablar de contenedor y entorno. El cuadro de entrada de un Chatbot es solo una puerta de acceso que entrega al modelo la intención del usuario. El entorno, en cambio, influye en cómo el usuario se expresa, actúa, da retroalimentación y comparte sus obras. El entorno es el reactor de las personas y también determina cómo se organizan las capacidades del modelo.

\begin{figure}[H]
\centering
\includegraphics[width=0.94\textwidth]{container-environment.png}
\caption{Cascarón, contenedor y entorno: ofrecer solo el cuadro de entrada de un Chatbot no basta; el entorno es el reactor de las personas. Diagrama conceptual propio, elaborado a partir de la conversación entre 01:16:33 y 01:21:18.}
\end{figure}

\begin{knowledgebox}{Lectura de la figura: Shell es solo la puerta de acceso; Environment es lo que moldea la conducta}
A la izquierda, Shell solo recibe la entrada y devuelve la salida; a la derecha, Environment aloja el contexto, las acciones, la retroalimentación y las reacciones de las personas. La hipótesis de producto de YouWare es que la creación con IA necesita un entorno donde las obras puedan publicarse, generar interacción, compartirse y seguir evolucionando, en lugar de entregarle el código al usuario para que lo despliegue por su cuenta.
\end{knowledgebox}

\subsection{Coding es como Camera: una herramienta nueva crea un grupo de personas nuevo}

Esta sección explica por qué YouWare apuesta por la creación mediante coding. Ming Chaoping compara coding con camera: en los primeros tiempos, hablar de cámaras era hablar de quienes usaban una réflex; con la llegada de la fotografía con celular surgió un grupo nuevo, el de los «fotógrafos con celular». Con coding ocurre lo mismo. Antes, coding era cosa de programadores; con AI coding aparecerá el Vibe Coder: alguien que no necesariamente domina todo el sistema de ingeniería, pero que puede generar obras con el modelo, su criterio estético y su intención.

\begin{figure}[H]
\centering
\includegraphics[width=0.96\textwidth]{vibe-coder-new-crowd.png}
\caption{El nuevo grupo de Vibe Coders: Coding es como Camera; el cambio de herramientas crea un grupo nuevo de creadores. Diagrama conceptual propio, elaborado a partir de la conversación entre 01:21:18 y 01:37:12.}
\end{figure}

\begin{knowledgebox}{Lectura de la figura: el grupo nuevo es una variable temprana de la tendencia}
En la figura, Camera pasa de los usuarios de réflex a los fotógrafos con celular, y Coding pasa de los ingenieros al Vibe Coder. Lo decisivo no es cuánto mejora la eficiencia del grupo profesional existente, sino si la herramienta nueva crea un grupo nuevo más grande y formas nuevas de contenido.
\end{knowledgebox}

\begin{importantbox}{El punto de entrada mínimo de YouWare}
La primera vez que Ming Chaoping escribió YouWare fue para resolver un problema: el HTML, los sitios web y los minijuegos escritos por IA no podían compartirlos ni usarlos las personas comunes. Pegar el código y obtener un sitio web accesible: ese punto de entrada saca coding del entorno local de ingeniería y lo lleva a un entorno de creación y distribución.
\end{importantbox}

\subsection{Resumen de la sección}

El contenedor y el entorno son el valor central de una empresa de aplicaciones de IA. Los modelos se volverán más potentes, pero el usuario necesita un espacio capaz de sostener la creación, el intercambio, la interacción y la retroalimentación. YouWare busca convertir coding de una tarea de ingeniería en una nueva forma de creación.
\section{El ecosistema de Agent: de AgentRank a OS Agent}

La sección anterior trató el contenedor de producto de YouWare; esta sección pasa a la valoración del ecosistema de Agent. Ming Chaoping (明超平) compara el ecosistema de Agent con la sociedad humana: puede haber un OS Agent que coordine a todos los Agent verticales, o bien un conjunto de Agent verticales que se invocan entre sí. En el futuro no importará solo si un Agent individual puede completar una tarea, sino quién puede ser invocado, en quién se confía y quién logra entrar en el efecto de red.

\begin{figure}[H]
\centering
\includegraphics[width=0.94\textwidth]{agent-ecosystem-two-models.png}
\caption{Dos ecosistemas de Agent: pequeño y fuerte al estilo de Singapur frente a grande y completo al estilo de Estados Unidos. Diagrama conceptual propio, elaborado a partir de la conversación entre 01:37:12 y 01:40:44.}
\end{figure}

\begin{knowledgebox}{Lectura de la figura: el ecosistema tiene dos formas de organización}
A la izquierda, Singapore-like representa lo pequeño y fuerte, con conexiones abiertas y dependencia del ecosistema externo; a la derecha, US-like representa lo grande y completo, autosuficiente y con capacidades de plataforma completas. Una empresa de aplicaciones debe decidir primero si se convertirá en un nodo de capacidad vertical o si buscará ser el coordinador a nivel de OS.
\end{knowledgebox}

\subsection{De PageRank a AgentRank}

La subsección anterior trató las formas del ecosistema de Agent; esta examina la clasificación y la distribución. En la era de la web, PageRank mide la importancia de una página mediante sus relaciones de enlace. En la era de los Agent, un mecanismo similar podría convertirse en AgentRank: las relaciones de invocación entre usuarios, OS Agent, Agent verticales, herramientas y contenidos formarán nuevos mecanismos de clasificación y distribución. Quien sea invocado por más nodos de alta confianza podrá obtener más tareas, datos e ingresos.

\begin{figure}[H]
\centering
\includegraphics[width=0.96\textwidth]{page-rank-to-agent-rank.png}
\caption{De PageRank a AgentRank: las páginas se clasifican por enlaces; los Agent podrían clasificarse por invocaciones y confianza. Diagrama conceptual propio, elaborado a partir de la conversación entre 01:40:44 y 01:44:02.}
\end{figure}

\begin{knowledgebox}{Lectura de la figura: los enlaces se vuelven invocaciones y el peso se vuelve confianza}
En la figura, Page, Links y PageRank se trasladan a Agent y AgentRank. La «clasificación» en el ecosistema de Agent no se limita a los resultados de búsqueda: combina la distribución de tareas, la invocación de herramientas, la autenticación de identidad y la confiabilidad de los resultados.
\end{knowledgebox}

\begin{knowledgebox}{Términos clave: la red de Agent}
\noindent\begin{tabularx}{\textwidth}{p{0.18\textwidth}p{0.36\textwidth}X}
\toprule
Término & Significado & Función en este episodio \\
\midrule
OS Agent & Agente de nivel superior que coordina a otros Agent y al entorno & Podría convertirse en el punto de entrada de la próxima generación. \\
Vertical Agent & Agente dedicado a una tarea vertical específica & Podría convertirse en un nodo de capacidad coordinado por otros. \\
AgentRank & Clasificación basada en invocaciones, confianza y resultados & Determina la distribución del tráfico y de las tareas. \\
Network Effect & Cuantos más usuarios y Agent entran, mayor es el valor & Una de las barreras de largo plazo de una empresa de aplicaciones. \\
Identity & Identidad, permisos y prueba de origen de un Agent & Resuelve problemas de confianza y seguridad. \\
\bottomrule
\end{tabularx}
\end{knowledgebox}

\subsection{Efecto de red y problema de confianza entre Agent}

Cuando los Agent se invocan entre sí, aparece el efecto de red, pero también aparece el problema de la confianza. Ming Chaoping menciona que en el futuro los Agent formarán tribus, como en la sociedad humana, y necesitarán documentos de identidad, cerraduras con contraseña y gestión de la confianza. Un Agent no puede invocar a otro Agent en nombre del usuario de forma arbitraria, ni entregar sin control los datos del usuario o ejecutar operaciones de alto riesgo. La identidad, los permisos, la auditoría y la confianza se convertirán en la infraestructura del ecosistema de Agent.

\begin{figure}[H]
\centering
\includegraphics[width=0.92\textwidth]{agent-network-effect.png}
\caption{Efecto de red de los Agent: las invocaciones, la identidad, la confianza y las herramientas entre Agent forman un ecosistema. Diagrama conceptual propio, elaborado a partir de la conversación entre 01:44:02 y 02:03:33.}
\end{figure}

\begin{knowledgebox}{Lectura de la figura: el efecto de red y las restricciones de seguridad aparecen al mismo tiempo}
En el centro de la figura está Agent Network, rodeado por Identity, Trust, Tools, Users, Calls y Rank. Cuanto más valiosa es la red, más necesita identidad y permisos; de lo contrario, las relaciones de invocación se convierten en superficie de ataque.
\end{knowledgebox}

\noindent\begin{tabularx}{\textwidth}{p{0.18\textwidth}p{0.36\textwidth}X}
\toprule
Problema de gobernanza & Por qué aparece & Posibles mecanismos de producto \\
\midrule
Identidad & El Agent debe actuar en nombre de un usuario o una organización & Cuentas, firmas, permisos, credenciales. \\
Confianza & Los Agent se invocan y se recomiendan entre sí & Historial de invocaciones, calificaciones, auditoría, prueba de origen. \\
Permisos & No todos los Agent pueden acceder a todos los datos & Mínimo privilegio, confirmación del usuario, aislamiento de alcance. \\
Responsabilidad & Tras un error hay que saber quién lo desencadenó y quién lo ejecutó & Registros, reproducción, flujos de trabajo rastreables. \\
Clasificación & Hay demasiados nodos de capacidad y se necesita un mecanismo de distribución & AgentRank, asignación de tareas, recomendación según el contexto. \\
\bottomrule
\end{tabularx}

\subsection{Los Agent de hoy son como quien acaba de tomar un palo para el fuego}

Esta subsección explica la metáfora del título. Ming Chaoping dice que los Agent de hoy se parecen a un gorila que acaba de tomar una piedra o un palo para el fuego: muy toscos, propensos a romper cosas, pero con una dirección de enorme alcance. La metáfora subraya dos hechos: por un lado, los Agent de hoy aún no son estables, y el uso de herramientas, el contexto, el entorno, la identidad y la confianza son todavía muy primitivos; por otro, en cuanto las herramientas y el entorno maduren, la forma de producir se reorganizará con rapidez.

\begin{figure}[H]
\centering
\includegraphics[width=0.94\textwidth]{ape-with-fire-stick.png}
\caption{La etapa actual de los Agent: como un gorila que acaba de tomar un palo para el fuego, con capacidades toscas pero una dirección de enorme alcance. Diagrama conceptual propio, elaborado a partir de la conversación entre 01:46:30 y 01:49:50.}
\end{figure}

\begin{knowledgebox}{Lectura de la figura: Primitive y Potential deben verse al mismo tiempo}
A la izquierda, Primitive significa tosco, con muchos intentos fallidos y propenso a romper cosas; a la derecha, Potential significa que la forma de producir se reorganiza cuando las herramientas mejoran. Al leer la figura no hay que limitarse a burlarse de la inestabilidad actual, pero tampoco hay que pasar por alto los límites actuales de las capacidades.
\end{knowledgebox}

\subsection{Resumen de la sección}

La oportunidad de largo plazo de las aplicaciones de Agent no está en herramientas aisladas, sino en el ecosistema, la clasificación, la identidad y el entorno. Lo que YouWare quiere construir no es solo una herramienta de coding, sino un entorno que pueda sostener la creación y la relación con el OS Agent.
\section{El modelo mejorará; las empresas de aplicaciones aportan el entorno y la experiencia}

La sección anterior trató el ecosistema de agentes; esta trata los límites de las empresas de aplicaciones. Ming Chaoping (明超平) dijo: «Confía siempre en que el modelo mejorará y en que el modelo no tiene nada que ver contigo». La frase no es pesimista; marca un límite para las empresas de aplicaciones: las empresas de modelos fundacionales crean personas más inteligentes, y las empresas de aplicaciones deben usar el entorno y la experiencia para que esos modelos más inteligentes se adapten a las necesidades reales de producción.

\begin{figure}[H]
\centering
\includegraphics[width=0.96\textwidth]{os-agent-living-system.png}
\caption{OS Agent: un sistema vivo. Los modelos fundacionales crean personas más inteligentes; las empresas de aplicaciones aportan el entorno y la experiencia. Diagrama conceptual propio, elaborado a partir de la conversación de 01:49:50--01:58:26.}
\end{figure}

\begin{knowledgebox}{Lectura de la figura: la empresa de aplicaciones no es una versión menor de la empresa de modelos}
En la figura, el Model se vuelve más inteligente, el Environment contiene los escenarios de uso, la Experience se adapta a las necesidades de producción, el OS Agent se convierte en un sistema vivo y el Feedback regresa de forma continua. El valor de la empresa de aplicaciones está en llevar la capacidad del modelo a flujos de trabajo repetibles y al entorno del usuario.
\end{knowledgebox}

\subsection{Exprimir la inteligencia con más eficiencia}

Ming Chaoping resume lo que observó en los últimos dos años como consumir tokens con más eficiencia y exprimir la inteligencia. En la era del Chatbot, el consumo de tokens dependía sobre todo del prompt del usuario; en la era del Agent, el sistema puede dividir las tareas por iniciativa propia, invocar herramientas y subagentes y generar productos intermedios, con lo que aumenta la velocidad de consumo de tokens. El problema es que ese consumo debe generar valor; de lo contrario, solo es costo.

\begin{importantbox}{La pregunta central para las empresas de aplicaciones}
No es «¿tengo un modelo?», sino «¿puedo convertir el modelo, con más eficiencia, en la capacidad de que el usuario complete sus tareas?». Esto exige que el producto entienda la tarea, el entorno, la retroalimentación, la distribución y el modelo de negocio, en lugar de solo envolver una API en otra capa.
\end{importantbox}

\subsection{Precios y modelo de negocio}

Antes se habló del entorno y del consumo de inteligencia; esta sección examina cómo se convierte en un modelo de negocio. En la entrevista se menciona que el precio de un producto Agent no tiene por qué basarse solo en el cobro por servicio. Pueden coexistir suscripciones, cobro por servicio, publicidad, incentivos para creadores y comisiones de plataforma. El juicio de Ming Chaoping es que el modelo de negocio debe generar beneficio mutuo para la plataforma y los creadores: que la empresa tenga utilidades suficientes para seguir desarrollando el producto y que los usuarios o creadores también reciban incentivos suficientes.

\begin{knowledgebox}{Términos clave: comercialización de aplicaciones de IA}
\noindent\begin{tabularx}{\textwidth}{p{0.20\textwidth}p{0.34\textwidth}X}
\toprule
Modelo & Escenario de uso & Riesgo \\
\midrule
Suscripción & Herramientas de uso frecuente, individuales o de equipo & Las empresas de modelos pueden incluirlas en sus paquetes y presionar el precio a la baja. \\
Cobro por servicio & Resultados de tarea definidos, como generar un sitio web, desplegarlo o ejecutar un flujo & Requiere controlar el costo y la calidad de la entrega. \\
Publicidad & Escenarios de contenido y distribución & Puede dañar la experiencia de creación. \\
Reparto con creadores & Ecosistema de la plataforma y comercio de obras & Requiere transacciones reales y gestión de derechos de autor. \\
Cobro a empresas & Colaboración en equipo, permisos, seguridad y despliegue & Ciclo de venta largo; el producto se vuelve pesado. \\
\bottomrule
\end{tabularx}
\end{knowledgebox}

\subsection{Por qué no basta con volverse to B}

Ming Chaoping mencionó que, si todas las empresas to C se convirtieran en empresas to B, habría resistencia generalizada. Detrás de esta frase está la cuestión de la ruta que sigue el emprendimiento en aplicaciones de IA. Muchos productos Agent se acercan a los servicios empresariales para tener ingresos seguros, pero si se vuelven completamente to B, el producto puede volverse pesado, el ciclo de venta se alarga, la comunidad de creación se debilita y se pierden los nuevos grupos de usuarios y los efectos de red.

Aquí está la tensión de YouWare: necesita comercializarse, pero no puede encerrarse demasiado pronto en el software empresarial tradicional. La creación con coding, la publicación de obras, los Vibe Coder y la red de Agent se parecen más a un ecosistema to C o prosumer; los permisos, la identidad, la colaboración en equipo y los pagos, en cambio, la empujan hacia el lado B. Una estrategia de producto exigente debería conservar las posibilidades de ambos lados, en lugar de sacrificar los accesos futuros a cambio de ingresos de corto plazo.

\begin{importantbox}{Elección de ruta: primero ser el entorno, después ser el producto que se vende}
Si YouWare se convierte primero en un entorno de creación, tendrá la oportunidad de acumular obras, usuarios, invocaciones de Agent y efectos de red; si demasiado pronto se dedica solo a entregas empresariales, puede obtener ingresos, pero le será difícil construir AgentRank y posicionarse en la mente de nuevos grupos de usuarios.
\end{importantbox}

\subsection{Resumen de la sección}

El valor de las empresas de aplicaciones no está en copiar los modelos fundacionales, sino en construir el entorno, la experiencia y la retroalimentación. La pregunta de largo plazo para YouWare es si puede pasar de ser un punto de entrada a la creación con coding a ser un sistema vivo que consuma inteligencia de forma continua, produzca obras y genere efectos de red.
\section{Marco para evaluar productos de Agent}

Las secciones anteriores trataron por separado la experiencia, la variable del modelo, el contenedor, AgentRank y OS Agent. En esta sección esas ideas se organizan en un marco reutilizable: al evaluar una aplicación de Agent no basta con preguntar «¿qué tan potente es el modelo?»; hay que preguntar si define un nuevo entorno de tareas, si crea un nuevo grupo de usuarios, si hace que los token produzcan valor de forma más eficiente, si genera efectos de red y si cuenta con mecanismos de identidad y confianza.

\begin{knowledgebox}{Nota de clase: un producto de Agent no es un script de automatización}
Un script de automatización termina al completar una tarea; un producto de Agent debe generar de forma continua tareas, retroalimentación, invocaciones y relaciones dentro de un entorno. La ambición de YouWare no es generar un sitio web de una sola vez, sino lograr que las obras hechas con AI coding puedan publicarse, recibir interacción, ser invocadas e integrarse en una nueva red de creación.
\end{knowledgebox}

\noindent\begin{tabularx}{\textwidth}{p{0.20\textwidth}p{0.34\textwidth}X}
\toprule
Dimensión de evaluación & Pregunta que hay que hacer & Correspondencia en YouWare \\
\midrule
Nuevo grupo de usuarios & ¿Aparecen personas que antes no hacían esto? & Un Vibe Coder no equivale a un programador tradicional. \\
Contenedor & ¿Las obras tienen un nuevo entorno donde alojarse y compartirse? & El HTML/code se convierte en un sitio web accesible e interactivo. \\
Eficiencia de la inteligencia & ¿Los token se convierten de forma eficiente en obras y valor? & Pasar de conversar con el prompt a generar productos utilizables. \\
Efectos de red & ¿Usuarios, obras y Agent se refuerzan entre sí? & La distribución de obras, las invocaciones y AgentRank forman un círculo virtuoso. \\
Gobernanza de la confianza & ¿Hay identidad, permisos, auditoría y límites de seguridad? & Cuando los Agent se organicen en tribus, necesitarán una credencial de identidad y un candado con contraseña. \\
Dependencia del modelo & Cuando el modelo mejora, ¿el producto se beneficia de forma natural? & Confiar siempre en que el Model mejorará, pero que la aplicación construya el entorno. \\
\bottomrule
\end{tabularx}

\begin{importantbox}{Una buena aplicación de Agent debe avanzar en la misma dirección en que mejora el modelo}
Si cada avance del modelo obliga a reescribir una capa de andamiaje del producto, el producto está mal posicionado; si con cada avance del modelo la calidad de las obras, el alcance de las tareas y el valor para el usuario dentro del entorno crecen de forma natural, el producto está apostando por la variable correcta.
\end{importantbox}

\subsection{Tres errores frecuentes}

Esta subsección mira el marco desde el lado opuesto. El primer error es convertir el producto de Agent en una ventana de chat. Una ventana de chat es lo bastante general, pero el entorno es demasiado delgado: el usuario no sabe cómo convertir la inteligencia en obras que pueda compartir. El segundo error es convertir el producto de Agent en software empresarial tradicional. El software empresarial tiene ingresos claros, pero tiende a volverse pesado y reduce el espacio para nuevos grupos de usuarios y para la red del ecosistema. El tercer error es tratar la capacidad del modelo como barrera de entrada de la aplicación. Las empresas de modelos seguirán mejorando; la barrera de la aplicación tiene que estar en el entorno del usuario, el retorno de datos, los efectos de red y los mecanismos de confianza.

\begin{warningbox}{La barrera de una aplicación no puede depender solo de lanzar primero la UI}
La UI actual es fácil de copiar. Lo realmente difícil de copiar son las relaciones con los usuarios, la red de obras, las relaciones de invocación entre Agent, la identidad y la confianza, y la retroalimentación continua. Si nada de esto se consolida, llegar primero solo da visibilidad, no una barrera.
\end{warningbox}

\subsection{Resumen de la sección}

El marco para evaluar productos de Agent se resume en una frase: observar si crea un nuevo entorno de tareas y un nuevo grupo de usuarios, y si hace que la capacidad del modelo siga generando valor dentro de ese entorno. La apuesta central de YouWare es que el AI coding pasará de ser una capacidad de ingeniería a ser una capacidad creativa, y que la capacidad creativa necesita nuevos contenedores y nuevas redes.
\section{Primera vez como CEO: valor emocional, Ego y mentalidad de jugador de ajedrez}

Esta sección pasa del producto a los mecanismos del fundador. Ming Chaoping (明超平) es CEO por primera vez: disfruta del financiamiento y de la atención, pero también siente que camina sobre hielo delgado. La ventana para las aplicaciones de IA es muy corta y la forma de los productos cambia con rapidez; el equipo necesita estabilidad emocional, y el fundador necesita combatir su Ego de manera constante. No se trata de un chisme sobre la personalidad de alguien, sino de cómo una organización emprendedora mantiene la calidad de su juicio en condiciones de incertidumbre extrema.

\begin{figure}[H]
\centering
\includegraphics[width=0.94\textwidth]{ego-countermeasure.png}
\caption{Combatir el Ego: quien es CEO por primera vez debe contenerse a sí mismo con una perspectiva de tercero, valor emocional y mentalidad de jugador de ajedrez. Diagrama conceptual propio, elaborado a partir de la conversación de 02:14:50--02:36:57.}
\end{figure}

\begin{knowledgebox}{Lectura de la figura: el mecanismo del CEO también es un mecanismo de producto}
La perspectiva de tercero proviene del debate; el valor emocional, de la gestión de la organización; caminar sobre hielo delgado, de la conciencia del riesgo; combatir el Ego, de la autocalibración; y la mentalidad de jugador de ajedrez, del análisis de las partidas y del juicio sobre el orden de las jugadas. En conjunto, determinan si el equipo puede seguir tomando las decisiones correctas.
\end{knowledgebox}

\subsection{El usuario por encima del Ego}

Ming Chaoping menciona una frase que circula en el equipo: el usuario está por encima del Ego. Cada quien tiene sus propias opiniones y también interpreta los datos a su manera; pero la verdadera calibración proviene de hablar con los usuarios, estar inmerso en Discord y observar cómo crean, cómo comparten y cómo comentan. Para una empresa de aplicaciones de IA, la retroalimentación de los usuarios es especialmente importante, porque los escenarios nuevos no tienen métricas consolidadas y muchas veces la intuición solo puede construirse a partir del uso real.

\begin{warningbox}{El Ego se disfraza de juicio estratégico}
La confianza del fundador es útil, pero también puede llevarlo a confundir su gusto personal, el éxito en el financiamiento y los elogios externos con necesidades de los usuarios. Que el usuario esté por encima del Ego no significa no tener opiniones, sino que las opiniones deben calibrarse continuamente con los escenarios de los usuarios.
\end{warningbox}

\subsection{El jugador y su contrincante}

La metáfora del «ajedrez» al final de la entrevista es fundamental. Ming Chaoping dice que lo importante no es reproducir el registro final de la partida, sino entender por qué cada jugada se hizo donde se hizo. Con los productos ocurre lo mismo: ver hoy los productos terminados de WeChat, Douyin, Instagram y TikTok no significa que copiar su forma baste para tener éxito. Lo decisivo son el orden, las prioridades y el momento: qué se hace ahora, qué se hace dentro de tres años y qué no se hace nunca.

\begin{importantbox}{Las prioridades del producto son el orden de las jugadas}
Si la jugada 50 es un error, quizá no haya jugada 100. Con el emprendimiento de aplicaciones de IA sucede igual: empezar por el contenedor de creación, la distribución y la red de Agents, o empezar por los servicios empresariales, las suscripciones y las capacidades del modelo, cambia todo el camino posterior.
\end{importantbox}

\subsection{Resumen de la sección}

La dificultad de ser CEO por primera vez no se reduce al financiamiento y la contratación: consiste en calibrarse a sí mismo de manera continua. La segunda mitad de EP101 muestra que emprender con aplicaciones de Agents exige que el juicio de producto, el estado emocional de la organización, la calibración con los usuarios y el combate contra el Ego operen al mismo tiempo.
\section{Síntesis y ampliación}

Esta sección condensa EP101 en seis juicios. Primero, el protagonismo en el emprendimiento de aplicaciones de IA se está desplazando de las empresas de modelos hacia las empresas de aplicaciones y los Agent, pero una empresa de aplicaciones no puede limitarse a ser un envoltorio de API. Segundo, la base de un gerente de producto proviene de la perspectiva de un tercero, la sensibilidad a la experiencia y la infraestructura de datos. Tercero, la bitter lesson exige que las empresas de aplicaciones sigan la dirección de las capacidades del modelo en lugar de ocultarlo con andamiajes antiguos. Cuarto, la velocidad de consumo de token y la per token valuation son nuevas herramientas de juicio de producto. Quinto, en el ecosistema de Agent surgirán AgentRank, efectos de red y problemas de identidad y confianza. Sexto, la autocalibración del CEO influye en la trayectoria del producto.

\begin{importantbox}{Ubicar EP101 en la serie de IA/internet de Zhang Xiaojun (张小珺)}
EP101 forma, junto con EP110/EP113/EP115, la línea temática de Agent: EP110 se orienta a los informes técnicos, EP113 al agentic LLM, EP115 a la filosofía de investigación de Agent, y EP101 es el terreno del emprendimiento de aplicaciones de IA: explica cómo las capacidades del modelo se convierten en producto mediante contenedores, entornos y efectos de red.
\end{importantbox}

\subsection{Tabla general de conceptos}

Esta sección reúne los conceptos de todo el texto. No es adorno terminológico, sino una forma de que el lector vea la cadena lógica de EP101: la base de producto determina cómo se identifica la experiencia; el contenedor y el entorno determinan cómo se usa el modelo; AgentRank y los efectos de red determinan cómo se forma el ecosistema; el usuario por encima del Ego determina cómo se calibra la organización de manera continua.

\noindent\begin{tabularx}{\textwidth}{p{0.18\textwidth}p{0.36\textwidth}X}
\toprule
Concepto & Significado & Función en este episodio \\
\midrule
Contenedor/entorno & Espacio de producto que aloja al modelo, a los usuarios, las obras y la retroalimentación & Sustituye la expresión despectiva de «envoltorio». \\
Vibe Coder & Nuevo grupo de usuarios que crea con AI coding & Análogo a lo que el fotógrafo con teléfono móvil es respecto de la cámara. \\
AgentRank & Mecanismo de invocación y clasificación por confianza de Agent & Explica la futura distribución de tareas y los efectos de red. \\
OS Agent & Agente de nivel superior que coordina entornos y Agent verticales & Podría convertirse en el punto de entrada de la próxima generación. \\
per token valuation & Valor para el usuario creado por cada token & Mide la eficiencia del consumo de inteligencia. \\
El usuario por encima del Ego & Calibrar las opiniones del fundador y del equipo con usuarios reales & Evita el ensimismamiento y las direcciones equivocadas. \\
\bottomrule
\end{tabularx}

\subsection{Preguntas para seguimiento}

\begin{enumerate}
\item ¿Podrá YouWare ampliarse desde el punto de entrada creativo de HTML-to-website hasta convertirse en un entorno de creación y distribución continuas?
\item ¿Los Vibe Coder son una novedad pasajera o formarán un nuevo grupo estable, como los fotógrafos con teléfono móvil?
\item ¿Surgirá realmente AgentRank, o será absorbido por los buscadores, las tiendas de aplicaciones y los puntos de entrada de modelos existentes?
\item ¿Quién se quedará con el OS Agent: las empresas de modelos, las de sistemas operativos, las de navegadores o las startups de aplicaciones?
\item ¿La identidad, los permisos, la confianza y la auditoría de los Agent se convertirán en un nuevo segmento de infraestructura?
\item ¿Cómo convertirá YouWare el entorno y la experiencia en una barrera defendible mientras los modelos siguen mejorando?
\end{enumerate}

\subsection{Comparación con Manus y Lovart}

Después de enumerar las preguntas para seguimiento, esta sección ubica EP101 en la trilogía de aplicaciones de Agent del primer semestre de 2025, con el fin de no ver todo emprendimiento de Agent como una misma forma. Manus se parece más a una narrativa de Agent de propósito general, que subraya que el mundo no es una extrapolación lineal y que hay que ser una variable dentro del juego estratégico; Lovart es un Agent vertical de diseño, que subraya los flujos de trabajo profesionales, la creación y los escenarios de los diseñadores; YouWare, en cambio, pone el foco en un nuevo contenedor para la creación mediante coding y en la posibilidad del OS Agent. Los tres muestran en conjunto que las aplicaciones de Agent no son un mismo producto: difieren en punto de entrada, entorno, usuarios, distribución y barreras.

\noindent\begin{tabularx}{\textwidth}{p{0.17\textwidth}p{0.30\textwidth}X}
\toprule
Caso & Posición central & Relación con YouWare \\
\midrule
Manus & Agent de propósito general; subraya la ejecución de tareas y la posición como variable & YouWare no copia directamente el Agent de propósito general, sino que busca un entorno de creación. \\
Lovart & Agent vertical de diseño; subraya los flujos de trabajo profesionales & YouWare también es vertical, pero su punto de entrada es la creación y publicación mediante coding. \\
YouWare & Contenedor de creación y potencial OS Agent & Se centra en el nuevo grupo de usuarios, el soporte de las obras, AgentRank y los efectos de red. \\
\bottomrule
\end{tabularx}

\begin{importantbox}{Conclusión común de la trilogía}
Lo central en las aplicaciones de Agent de 2025 no es «quién empaqueta primero una ventana de chat», sino quién logra encontrar nuevos entornos de tareas, nuevos grupos de usuarios y nuevos mecanismos de distribución. Manus, Lovart y YouWare apuestan cada uno por un nivel distinto: tareas de propósito general, diseño profesional y contenedor de creación mediante coding.
\end{importantbox}

\subsection{Lecturas adicionales}

\begin{itemize}
\item Si interesan los informes técnicos de Agent y la ingeniería de contexto, puede contrastarse con las notas de EP110 sobre los informes técnicos de Kimi K2 / ChatGPT Agent / Qwen3-Coder.
\item Si interesan el Agentic LLM y los límites de las empresas de modelos, puede contrastarse con la entrevista de EP113 a Yang Zhilin (杨植麟) sobre Kimi K2.
\item Si interesan la filosofía de investigación de Agent, la interface y el reward, puede contrastarse con la entrevista de EP115 a Yao Shunyu (姚顺雨).
\item Si interesa la ventana de oportunidad para emprender aplicaciones de IA, puede contrastarse con las entrevistas a Lovart / Chen Mian (陈冕) y a Manus / Xiao Hong (肖弘).
\end{itemize}

\end{document}
